\documentclass[12pt]{article}
\usepackage[utf8]{inputenc}
\usepackage[T1]{fontenc}
\usepackage[margin=1in]{geometry}
\usepackage{setspace}
\usepackage{xcolor}
\usepackage[colorlinks=true, urlcolor=blue!60!black, linkcolor=blue!60!black]{hyperref}
\setlength{\parindent}{2em}
\setlength{\parskip}{1.25em}
\renewcommand{\baselinestretch}{1.0}
\begin{document}
\begin{center}
{\LARGE \textit{Everson v. Board of Education}, 330 U.S. 1 (1947)\par}
\vspace{0.8em}
{\large Case Brief --- Four-Page Edition\par}
\end{center}
\vspace{1em}

\section{Memory Jogger}

Incorporating the Establishment Clause against the states and invoking Jefferson's ``wall of separation between church and State,'' the Court nonetheless upheld 5--4 New Jersey's reimbursement of bus fares for children attending parochial schools as a neutral public-welfare benefit, like police and fire protection (\href{https://supreme.justia.com/cases/federal/us/330/1/}{330 U.S. 1}).

\section{Detailed Case Facts}

A 1941 New Jersey statute let local school boards provide transportation for children, including to schools ``other than a public school, except such school as is operated for profit.'' Ewing Township, near Trenton, which had no public high school, reimbursed parents from tax funds for fares on regular public buses to three public high schools outside the district and four Catholic parochial schools. The parochial schools met state educational requirements but also gave regular Catholic religious instruction. Reimbursements averaged roughly \$40 per child a year.

\section{Procedural History}

Arch R. Everson, a township taxpayer, challenged the statute and resolution as applied to parochial-school students. The New Jersey Supreme Court, then an intermediate court, held the reimbursements unauthorized under the state constitution; the Court of Errors and Appeals reversed, finding no state or federal violation. Everson appealed the federal questions. He did not challenge the exclusion of for-profit private schools, and the record showed no township children attending non-Catholic private schools.

\section{Judicial Votes}

Affirmed, \textbf{5--4}, decided February 10, 1947.
\begin{itemize}
\item \textbf{Majority:} Justice Black, joined by Chief Justice Vinson and Justices Reed, Douglas, and Murphy.
\item \textbf{Concurrences:} None.
\item \textbf{Dissents:} Justice Jackson, joined by Justice Frankfurter; Justice Rutledge, joined by Justices Frankfurter, Jackson, and Burton.
\end{itemize}
All nine Justices agreed that the Establishment Clause is incorporated and demands strict separation.

\section{Holding}

(1) Reimbursing bus fares to accredited schools serves a public purpose and does not violate due process. (2) The Establishment Clause applies to the states through the Fourteenth Amendment with the same breadth as against the Federal Government. (3) Although the Clause erects a ``high and impregnable'' wall forbidding aid to one or all religions or any tax to support religious institutions, New Jersey did not breach it: a neutral, generally available welfare benefit to parents, like police and fire protection, gave no money to the schools.

\section{Analysis of Opinions}

\textbf{Black (majority).} Public-purpose review demands ``the most extreme caution,'' and under \textit{Cochran} it was ``much too late'' to deny that facilitating secular education serves a public purpose. Incorporating the Establishment Clause, Black read it through Virginia's disestablishment struggle, Madison's \emph{Memorial and Remonstrance} and Jefferson's \emph{Bill for Religious Liberty}, and listed strict prohibitions. He then pivoted: the Free Exercise Clause forbids excluding believers ``because of their faith, or lack of it, from receiving the benefits of public welfare legislation.'' Traffic police, fire protection, sewers, and sidewalks are ``so separate and so indisputably marked off from the religious function'' that withholding them would make the state religion's adversary, and under \textit{Pierce} parents may satisfy compulsory education at religious schools.

\textbf{Jackson (dissenting).} The opinion's separationist undertones are ``utterly discordant with its conclusion,'' like Byron's Julia, who protested and then consented. Ewing ran no buses and provided no safety service; it reimbursed only public and Catholic school students, applying a religious test. Citing Canon Law, he called parochial education the rock on which the Church rests, and warned that government regulates what it subsidizes.

\textbf{Rutledge (dissenting).} Appending Madison's \emph{Memorial and Remonstrance}, he argued that one word, ``religion,'' governs both clauses, so the Establishment Clause forbids ``every form of public aid or support for religion.'' Madison opposed even nondiscriminatory assessments, and transportation is as essential to religious schooling as books or salaries, so no rational line remains.

\textbf{Later developments.} \textit{McCollum} and \textit{Engel} enforced the wall while \textit{Allen} lent textbooks and \textit{Lemon} formalized a test; since \textit{Zelman} (2002), \textit{Trinity Lutheran}, \textit{Espinoza}, and \textit{Carson} turned \textit{Everson}'s neutrality strand into a mandate.

\section{Examples: Future Applications}

\begin{enumerate}
\item \textbf{Same-Side (Upheld): Free Transit Passes.} A city gives every K--12 student a free bus and subway pass, including religious-school students. The benefit goes to families, not schools, and is indistinguishable from \textit{Everson}'s fares.
\item \textbf{Same-Side (Upheld): Textbook Loans.} A state lends secular textbooks to all students, including at parochial schools, as \textit{Board of Education v.\ Allen} upheld.
\item \textbf{Opposite-Side (Invalid): Clergy Salaries.} A county grants equal tax funds to all faiths to pay clergy. Aid to all religions is forbidden, and the grants fund the core religious function.
\item \textbf{Opposite-Side (Invalid): Religious Classes in Public Schools.} Clergy teach weekly classes in public classrooms during school hours, as struck down in \textit{McCollum}.
\item \textbf{Fence-Sitter: Tuition Aid and Charter Schools.} A State excludes ``sectarian'' schools from tuition aid and requires charter schools to be nonsectarian. \textit{Carson} condemns the first; the second turns on whether charter schools are state actors.
\end{enumerate}

\section{Critique}

\textbf{Federalism and history.} Critics contend the original Establishment Clause barred federal interference with state establishments, making casual incorporation problematic, and fault Black and Rutledge for treating Virginia's 1785--1786 struggle as the sole gloss on 1789.

\textbf{Fungibility and regulatory strings.} As Rutledge noted, reimbursements free up funds for religious instruction; conversely, Jackson's warning that government regulates what it subsidizes now drives post-\textit{Carson} disputes over conditions on public aid.

\section{Key Quotations}

\begin{quote}
``The `establishment of religion' clause of the First Amendment means at least this: neither a state nor the Federal Government can set up a church. Neither can pass laws which aid one religion, aid all religions, or prefer one religion over another. Neither can force nor influence a person to go to or to remain away from church against his will or force him to profess a belief or disbelief in any religion. No person can be punished for entertaining or professing religious beliefs or disbeliefs, for church attendance or non-attendance. No tax in any amount, large or small, can be levied to support any religious activities or institutions, whatever they may be called, or whatever form they may adopt to teach or practice religion. Neither a state nor the Federal Government can, openly or secretly, participate in the affairs of any religious organizations or groups, and vice versa. In the words of Jefferson, the clause against establishment of religion by law was intended to erect `a wall of separation between church and State.'\,''
\end{quote}
\noindent --- Justice Black, Opinion of the Court, 330 U.S. at 15--16.

\begin{quote}
``That Amendment requires the state to be a neutral in its relations with groups of religious believers and nonbelievers; it does not require the state to be their adversary. State power is no more to be used so as to handicap religions than it is to favor them. \ldots{} The First Amendment has erected a wall between church and state. That wall must be kept high and impregnable. We could not approve the slightest breach. New Jersey has not breached it here.''
\end{quote}
\noindent --- Justice Black, Opinion of the Court, 330 U.S. at 18.

\begin{quote}
``But we cannot have it both ways. Religious teaching cannot be a private affair when the state seeks to impose regulations which infringe on it indirectly, and a public affair when it comes to taxing citizens of one faith to aid another, or those of no faith to aid all. \ldots{} If the state may aid these religious schools, it may therefore regulate them. Many groups have sought aid from tax funds, only to find that it carried political controls with it.''
\end{quote}
\noindent --- Justice Jackson, dissenting, 330 U.S. at 27--28.

\end{document}
